\section{Introducción}

\begin{frame}
\frametitle{Contexto}
\begin{itemize}
    \item El análisis de componentes principales (PCA) es una técnica matemática utilizada para problemas de reducción de dimensionalidad,
    compresión de datos con pérdidas, extracción de componentes y visualización de datos.
    \item En esencia, es un problema de optimización, definida de dos posibles manera, resultando en un mismo algoritmo.
\end{itemize}
\end{frame}


\begin{frame}
\frametitle{Contexto}
\begin{itemize}
    \item El primero corresponde a una proyección ortogonal de los datos en un espación de lineal de menor dimensión, en el cual la \textbf{varianza} es maximizada
    \item El segundo corresponde a una proyección lineal que \textbf{minimiza} el costo de proyección promedio, definido como la \text{mean squared distance} entre los datos y sus proyecciones.
\end{itemize}
\end{frame}

\begin{frame}
\frametitle{Contexto}
\begin{itemize}
    \item Para aplicar el criterio de máxima varianza, se busca obtener el conjunto de vectores unitarios ordenados por la varianza proyectada sobre ello.
    \item Para ello, se calcula la \textbf{matriz de covarianza}, una matriz cuadrada de tamaño $n \times n$, con $n$ el número de valores de la matriz original.
\end{itemize}
\end{frame}

\begin{frame}
\frametitle{Problemas actuales}
\begin{itemize}
    \item En el marco de implementar un algoritmo de PCA computacionalmente, el cálculo de la matriz de covarianza resulta costoso en materia de memoria y tiempo de ejecución para matrices muy grandes.
    \item Por lo tanto, se utilizan técnicas alternativas para estimar los componentes principales de la matriz de origen.
\end{itemize}
\end{frame}

\begin{frame}
\frametitle{Problemas actuales}
\begin{itemize}
    \item Entre las técnicas alternativas se encuentra la \textit{descomposición de valores singulares} (SVD), que busca descomponer la matriz $X$ de origen en tres submatrices, específicas:
    \begin{equation}
        X = U\Sigma V^T
    \end{equation}
    \item $\Sigma$ corresponde a los \textbf{valores singulares} de la matriz $X$ mientras que $U$ y $V$ contienen información sobre los espacios que basan a las filas y columnas de la matriz.
\end{itemize}
\end{frame}

\begin{frame}
\frametitle{Problemas actuales}
\begin{itemize}
    \item El algoritmo de SVD requiere realizar un proceso de ortogonalización entre pares de columnas de manera iterativa, con un valor de tolerancia $\epsilon$.
    \item Para matrices muy grandes, este proceso realizado secuencialmente puede resultar en una complejidad cercana a $\mathcal{O}(n^{1.5})$, para una matriz de tamaño $T \times D$ y $n = TD$.
\end{itemize}    
\end{frame}

\begin{frame}
\frametitle{Motivación}
\begin{itemize}
    \item Se propone dentro del marco del curso, analizar los efectos de implementar el algoritmo en paralelo usando CUDA, en términos de complejidad temporal.
\end{itemize}
\end{frame}
